% ======================================================================
% Route r17-oracle-barrier.  Paste-ready; the full file is also at
% <code_dir>/result.tex .  Labels prefixed ob:.
% ======================================================================

\section{One model for the propagation mechanisms, and a family every
algorithm of it loses on}\label{sec:ob}

The mechanisms $M1$--$M6$ of \Cref{sec:gap} are presented as six devices
with six barriers. They are one device. What separates $M1$ from $M2$ is
not that one is arithmetic-free and the other is not: it is the
\emph{grade set} over which the same clauses are evaluated, together with
the fact that the coarse one is read coinductively and the fine one
inductively (\Cref{ob:m1-grade}). Once that is said, all six, with the
transport programme, the min-plus closure, and every combination and
mutual seeding of them, fit into one model with one cost measure, and one
certificate bounds all of them at once (\Cref{ob:invariance}). This
section defines the model, proves the membership statements, proves the
certificate theorem, and exhibits a one-player family on which every
algorithm of the model is silent at every controlled vertex for
$2^{\Omega(N)}$ steps --- including the composite that runs $M1$ first and
seeds the rest with its verdicts, which is what disposes of every family
recorded before.

\subsection{The model}

Throughout $G$ is an SSG on $V$,
$\mathcal X=\{z\in[0,1]^{V}:z(t_0)=0,\ z(t_1)=1\}$, $Q(G;L,U)$ is the
transport polytope of \Cref{def:transport}, and $Z_0,Z_1$ are the sets of
\Cref{def:simorder}, computable arithmetic-free by
\Cref{prop:qualitative}. The \emph{free bracket} is
$L_Z:=\mathbf 1_{Z_1}$, $U_Z:=\mathbf 1-\mathbf 1_{Z_0}$. For nonempty
compact $P\subseteq\mathcal X$ put
\[
  \Lambda(P)(x,y):=\max_{z\in P}\bigl(z(x)-z(y)\bigr),\qquad
  \Lambda_{\times}(P)(x,y):=
  \begin{cases}
    0,&z(x)=0\text{ for all }z\in P,\\
    +\infty,&z(y)=0<z(x)\text{ for some }z\in P,\\
    \max\{z(x)/z(y):z\in P,\ z(y)>0\},&\text{otherwise,}
  \end{cases}
\]
and $K(\Delta):=\{z\in Q(G;L_Z,U_Z):z(x)-z(y)\le\Delta(x,y)\ \forall x,y\}$.

\begin{definition}[The difference-propagation model $\mathfrak D$]\label{ob:model}
A \emph{state} is a pair $(\Delta,R)$ with
$\Delta\colon V\times V\to[-1,1]$ and $R\colon V\times V\to[0,+\infty]$;
it is \emph{sound} on $G$ if $\Delta(x,y)\ge w^\ast(x)-w^\ast(y)$ and
$w^\ast(x)\le R(x,y)\,w^\ast(y)$ for all $x,y$. The \emph{repertoire}
$\mathcal R$ consists of the following operators, each leaving the other
register unchanged unless stated.
\begin{enumerate}[label=\textup{(D\arabic*)},leftmargin=3.4em]
\item $\Delta\mapsto\min(\Delta,\mathcal L\Delta)$, where $\mathcal L$ is
  one round of the clauses of \Cref{def:slack};
\item $\Delta\mapsto\mathcal C\Delta$, the min-plus closure of
  \Cref{def:trans-slack};
\item $\Delta\mapsto\min\bigl(\Delta,\Lambda(K(\Delta))\bigr)$, the
  transport step of \Cref{def:transport} with the free bracket;
\item $R\mapsto\min(R,\mathcal R_{\times}R)$, where $\mathcal R_{\times}$
  is one round of the clauses of \Cref{def:ratio};
\item $\Delta\mapsto\rho_{\Gamma}\Delta$ and $R\mapsto\rho_{\Gamma}R$,
  where the \emph{grade set} $\Gamma$ is a finite subset of $[-1,1]$
  containing $1$, resp.\ of $[0,+\infty]$ containing $+\infty$, and
  $\rho_\Gamma$ rounds every entry \emph{up} to the least member of
  $\Gamma$ that is at least as large;
\item $(\Delta,R)\mapsto(\min(\Delta,\Delta'),\min(R,R'))$, where
  $(\Delta',R')$ is a state reached by another run of the model on the
  same $G$ (combination, and mutual seeding);
\item $\Delta(x,y)\leftarrow0$ wherever $R(x,y)\le1$, and
  $R(x,y)\leftarrow1$ wherever $\Delta(x,y)\le0$ (order transfer between
  the registers);
\item replacement of the state by the limit of iterating any composite
  $\Theta$ of \textup{(D1)}--\textup{(D7)} whose grade sets are finite,
  from the current state or from the constant least state; the limit
  exists because $\Theta$ is monotone and the grade sets are finite.
\end{enumerate}
A member of $\mathfrak D$ is an algorithm which, on every stopping $G$,
starts from
$\Delta_0(x,y)=\mathrm{clamp}_{[-1,1]}(U_Z(x)-L_Z(y))$ and
$R_0\equiv+\infty$, applies operators of $\mathcal R$, is sound whenever
it reports, and reports as in \Cref{rem:own-successor}: at $v\in\Vmax$ it
may return $(v,i)$ when $\Delta(v,v^{(i)})\le0$ or $R(v,v^{(i)})\le1$
(clause~(i)), when $\Delta(v^{(1-i)},v)<0$ or $R(v^{(1-i)},v)<1$
(clause~(ii)), or by the pair tests $\Delta(v^{(1-i)},v^{(i)})\le0$,
$R(v^{(1-i)},v^{(i)})\le1$; dually at $\Vmin$. Its \emph{cost} on $G$ is
the number of operator applications made before it reports, counted over
the whole computation, including the runs invoked by \textup{(D6)}. Write
$\mathfrak D_0$ for the members using no \textup{(D8)} step.
\end{definition}

\begin{theorem}[$M1$ is the slack calculus over the two-element grade
set]\label{ob:m1-grade}
Let $G$ be an SSG, let $\mathcal L$ be the clause set of
\Cref{def:slack}, and let $\mathcal F'$ be the operator of
\Cref{def:simorder} with the diagonal adjoined to its base --- a local
matching calculus in the sense of \Cref{def:lmc}, which allows exactly
that. For $R\subseteq V\times V$ write $\chi_R(x,y):=0$ if $(x,y)\in R$
and $1$ otherwise. Then
\[
  \mathcal F'(R)\;=\;\{(x,y):(\mathcal L\chi_R)(x,y)\le0\},
  \qquad\text{that is}\qquad
  \chi_{\mathcal F'(R)}\;=\;\rho_{\{0,1\}}\,\mathcal L\,\chi_R ,
\]
and $\operatorname{gfp}(\mathcal F')=\operatorname{gfp}(\mathcal F)=\preceq$.
Consequently $\preceq$ is produced by one \textup{(D8)} step with
$\Gamma=\{0,1\}$ and $\Theta=\textup{(D1)}$ from the constant least state:
\Cref{def:simorder} \emph{is} \Cref{def:slack}, read over the grade set
$\{0,1\}$ and coinductively instead of over the dyadic grades and
inductively.
\end{theorem}

\begin{proof}
Entrywise. $\mathcal L\chi_R$ is the minimum, over the applicable clauses,
of quantities each of which is a $\min$, a $\max$ or a mean of entries of
$\chi_R\in\{0,1\}$, together with the constants $1$; $0$ when $x\in Z_0$,
when $y\in Z_1$ or when $x=y$; and $-1$ when $x\in Z_0$ and $y\in Z_1$.
Such a quantity is $\le0$ exactly when it is $0$ or $-1$; and a mean of
two $\{0,1\}$ entries is $0$ iff both are.

The constants give the base $(Z_0\times V)\cup(V\times Z_1)$ together with
the diagonal, which is the base of $\mathcal F'$.

(down) at a non-sink $y$: $\mathrm{ag}^{-}_{y}$ is $\min$ at $\Vmax$,
hence $0$ iff $(x,y^{(j)})\in R$ for \emph{some} $j$; it is $\max$ at
$\Vmin$ and the mean at $\Vavg$, hence $0$ iff for \emph{both} $j$. That
is clause~(i) of \Cref{def:simorder}. (up) at a non-sink $x$:
$\mathrm{ag}^{+}_{x}$ is $\min$ at $\Vmin$ ($0$ iff for some $i$), $\max$
at $\Vmax$ and the mean at $\Vavg$ ($0$ iff for both $i$); clause~(ii).
Composite, $x,y$ non-sinks of the same type: at $\Vmax$ the clause is
$\max_i\min_j\chi_R(x^{(i)},y^{(j)})$, which is $0$ iff every $i$ admits a
$j$; at $\Vmin$ it is $\max_j\min_i$, $0$ iff every $j$ admits an $i$; at
$\Vavg$ it is
$\tfrac12(\chi_R(x^{(0)},y^{(\pi(0))})+\chi_R(x^{(1)},y^{(\pi(1))}))$ for
either bijection $\pi$, and the minimum over the two bijections is $0$ iff
some bijection has both entries $0$. That is clause~(iii). No clause
occurs on one side only, so the two sets agree.

$\operatorname{gfp}(\mathcal F')=\operatorname{gfp}(\mathcal F)$:
$\mathcal F'\supseteq\mathcal F$ gives $\supseteq$. Conversely let
$R=\mathcal F'(R)=\mathcal F(R)\cup\Delta_{\mathrm{diag}}$. The diagonal
is a post-fixed point of $\mathcal F$ (the proof of
\Cref{thm:simorder-sound}), so
$\Delta_{\mathrm{diag}}\subseteq\mathcal F(\Delta_{\mathrm{diag}})
\subseteq\mathcal F(R)$ by monotonicity, whence $R\subseteq\mathcal F(R)$
and $R\subseteq\operatorname{gfp}(\mathcal F)$.

Finally $\operatorname{gfp}(\mathcal F')=\bigcap_k\mathcal F'^{\,k}(V\times V)$
in the finite lattice $2^{V\times V}$, and $\chi_{V\times V}\equiv0$ is
the least element of $\{0,1\}^{V\times V}$, so by the displayed identity
the sequence $\chi_{\mathcal F'^{\,k}(V\times V)}$ is exactly the
iteration of $\rho_{\{0,1\}}\mathcal L$ from the least state, whose limit
\textup{(D8)} returns.
\end{proof}

\begin{lemma}[The mechanisms are members]\label{ob:mem}
On every stopping $G$, with $k$ rounds of the mechanism costing $O(k)$
operator applications:
\textup{(a)} the slack calculus $M2$ of \Cref{def:slack} is
\textup{(D1)}$^{k}$;
\textup{(b)} the transitive slack calculus $M2T$ of
\Cref{def:trans-slack} is $(\textup{(D1)}\,\textup{(D2)})^{k}$;
\textup{(c)} the transport test $M4$ of \Cref{def:transport} with the free
seed, read by \Cref{rem:own-successor}, is one \textup{(D3)};
\textup{(d)} the hybrid $M5$ of \Cref{def:hybrid} is
$(\textup{(D1)}\,\textup{(D2)}\,\textup{(D3)}\,\textup{(D2)})^{k}$;
\textup{(e)} the ratio calculus $M6$ of \Cref{def:ratio} is
\textup{(D4)}$^{k}$;
\textup{(f)} the preorder $M1$ of \Cref{def:simorder} is one
\textup{(D8)} step, by \Cref{ob:m1-grade};
\textup{(g)} every combination of two members, and every composite that
seeds one with the output of another, is a member, by \textup{(D6)} and
\textup{(D7)}.
No member of $\mathfrak D$ evaluates a strategy: $\mathcal R$ contains no
operator whose output depends on the inverse of $I-A_{\sigma,\tau}$ for a
positional pair. Accordingly $M3=M(p,k)$ of \Cref{def:seeded},
all-switches, \Cref{def:bsi} and $R_{\mathrm{BR}}$ are \emph{not}
members; write $\mathfrak D^{+}$ for $\mathfrak D$ together with
\begin{enumerate}[label=\textup{(E)},leftmargin=3.4em]
\item $\Delta\mapsto\min\bigl(\Delta,\ (x,y)\mapsto U(x)-L(y)\bigr)$ for
  $L=\val_\sigma$, $U=\val^{\tau}$ and any positional pair
  $(\sigma,\tau)$ --- one exact policy evaluation
  (\Cref{thm:eval-stopfree}),
\end{enumerate}
which contains all of them.
\end{lemma}

\begin{proof}
(a),(b),(d),(e) transcribe the definitions; the running minimum with the
previous matrix in \Cref{def:slack,def:ratio} is built into
\textup{(D1)},\textup{(D4)}, and \Cref{lem:hybrid-cutting} shows that in
\Cref{def:hybrid} the outer minimum and the outer closure are redundant,
which is why the schedule in (d) is the displayed one. (c) is
\Cref{rem:own-successor} together with
$\operatorname{Sep}(p,q)=\Lambda(K(\mathbf1))(q,p)$. (f) is
\Cref{ob:m1-grade}. (g): a componentwise minimum of sound states is
sound, and \textup{(D7)} is sound because $R(x,y)\le1$ and
$\Delta(x,y)\le0$ both say $w^\ast(x)\le w^\ast(y)$. The last sentence
reads \Cref{def:seeded} and \Cref{def:missing}: the seed of $M(p,k)$ is
$\val_{\sigma_p}$, and all-switches, \Cref{def:bsi} and
$R_{\mathrm{BR}}$ evaluate $\val_\sigma$ (and $\val^{\tau}$) exactly at
every round.
\end{proof}

\begin{remark}[Where the model has to stop]\label{ob:scope}
Excluding \textup{(E)} is forced if a barrier is to be provable at all.
\Cref{thm:seed-dichotomy} says that an instance on which $M(p,0)$ is
silent is one on which all-switches has not converged after $p$ rounds; a
barrier for $\mathfrak D^{+}$ with $p$ polynomial would therefore be a
superpolynomial all-switches family, which this document does not have
(\Cref{rem:blowup-realise}). A barrier for $\mathfrak D$ says exactly
this: no amount of \emph{propagation} --- order, additive, multiplicative,
convex, coarse or fine, coinductive or inductive, in any combination and
any order --- replaces one exact solve of a linear system.
\end{remark}

\subsection{Certificates bound the whole model}

\begin{lemma}\label{ob:sign}
For nonempty compact $P\subseteq\mathcal X$ and all $x,y$:
$\Lambda_{\times}(P)(x,y)\le1$ if and only if $\Lambda(P)(x,y)\le0$.
\end{lemma}

\begin{proof}
If $\Lambda(P)(x,y)\le0$ then $z(x)\le z(y)$ for every $z\in P$; in
particular $z(y)=0$ forces $z(x)=0$, so the $+\infty$ case does not occur
and every ratio with positive denominator is at most $1$. Conversely
$\Lambda_\times(P)(x,y)\le1$ excludes the $+\infty$ case, so $z(y)=0$
implies $z(x)=0$, and where $z(y)>0$ the ratio bound gives
$z(x)\le z(y)$.
\end{proof}

\begin{theorem}[Convex certificates bound every member of
$\mathfrak D_0$]\label{ob:invariance}
Let $P_0\supseteq P_1\supseteq\cdots\supseteq P_K$ be nonempty compact
convex subsets of $Q(G;L_Z,U_Z)$ such that for every $k<K$
\begin{enumerate}[label=\textup{(\roman*)},leftmargin=2.6em]
\item every clause of \Cref{def:slack} evaluated at $\Lambda(P_k)$ is at
  least $\Lambda(P_{k+1})$ entrywise, and
\item every clause of \Cref{def:ratio} evaluated at
  $\Lambda_{\times}(P_k)$ is at least $\Lambda_{\times}(P_{k+1})$
  entrywise.
\end{enumerate}
Then every member of $\mathfrak D_0$ has, after $k\le K$ operator
applications, a state with $\Delta\ge\Lambda(P_k)$ and
$R\ge\Lambda_{\times}(P_k)$. Consequently, within $K$ applications it
decides $v\in\Vmax$ only if one of the four quantities
$\Lambda(P_K)(v,v^{(0)})$, $\Lambda(P_K)(v,v^{(1)})$,
$\Lambda(P_K)(v^{(0)},v^{(1)})$, $\Lambda(P_K)(v^{(1)},v^{(0)})$ is
$\le 0$, or one of $\Lambda(P_K)(v^{(0)},v)$, $\Lambda(P_K)(v^{(1)},v)$
is $<0$, or the corresponding multiplicative quantity is $\le1$
resp.\ $<1$; dually at $\Vmin$.
\end{theorem}

\begin{proof}
Induction on the number $k$ of applications, with invariant
$\Delta\ge\Lambda(P_k)$ and $R\ge\Lambda_{\times}(P_k)$; both sequences
are nonincreasing in $k$ because the $P_k$ are nested.

\emph{Start.} Every $z\in P_0\subseteq Q(G;L_Z,U_Z)$ has
$L_Z\le z\le U_Z$, so $z(x)-z(y)\le U_Z(x)-L_Z(y)$, and $\Lambda$ is
$\ge-1$ on $\mathcal X$, so $\Lambda(P_0)\le\Delta_0$;
$R_0\equiv+\infty\ge\Lambda_\times(P_0)$.

\textup{(D1)}: the clauses are monotone in the matrix, so
$\mathcal L\Delta\ge\mathcal L\Lambda(P_k)\ge\Lambda(P_{k+1})$ by (i), and
$\Delta\ge\Lambda(P_k)\ge\Lambda(P_{k+1})$.
\textup{(D2)}: $\Lambda(P_{k+1})$ obeys the triangle inequality
(\Cref{lem:hybrid-cutting}), so the min-plus closure of a matrix
$\ge\Lambda(P_{k+1})$ stays $\ge\Lambda(P_{k+1})$.
\textup{(D3)}: every $z\in P_{k+1}$ lies in $Q(G;L_Z,U_Z)$ and satisfies
$z(x)-z(y)\le\Lambda(P_{k+1})(x,y)\le\Delta(x,y)$, so
$P_{k+1}\subseteq K(\Delta)$ and $\Lambda(K(\Delta))\ge\Lambda(P_{k+1})$.
\textup{(D4)}: as \textup{(D1)}, with (ii).
\textup{(D5)}: $\rho_\Gamma$ only raises entries.
\textup{(D6)}: the cost counts the whole computation, so the run
supplying $(\Delta',R')$ has made at most $k$ applications and
$\Delta'\ge\Lambda(P_k)\ge\Lambda(P_{k+1})$, likewise for $R'$; a minimum
of two such matrices is again $\ge\Lambda(P_{k+1})$.
\textup{(D7)}: if $R(x,y)\le1$ then $\Lambda_\times(P_k)(x,y)\le1$, so
$\Lambda(P_k)(x,y)\le0$ by \Cref{ob:sign} and setting $\Delta(x,y):=0$
preserves the invariant at that entry; symmetrically in the other
direction.

The last sentence is the reporting rule of \Cref{ob:model} applied to a
state satisfying the invariant.
\end{proof}

\begin{proposition}[Every barrier for $\mathfrak D$ is a rate
barrier]\label{ob:no-permanent}
Let $G$ be stopping and $P\subseteq Q(G;L_Z,U_Z)$ nonempty compact convex
with $\mathcal L\Lambda(P)\ge\Lambda(P)$ entrywise (a \emph{stationary}
certificate). Then $\Lambda(P)(x,y)\le w^\ast(x)-w^\ast(y)$ for all
$x,y$; in particular $\Lambda(P)(v,v^{(i)})\le0$ at every $v\in\Vmax$ and
every optimal successor $v^{(i)}$. So no stationary certificate blocks
clause~(i) of \Cref{rem:own-successor}, and every lower bound obtainable
from \Cref{ob:invariance} is valid for finitely many rounds only.
\end{proposition}

\begin{proof}
Take the constant chain $P_k:=P$; hypothesis (i) of
\Cref{ob:invariance} is the assumption, so the member \textup{(D1)}$^k$
--- the slack calculus --- satisfies $\Delta_k\ge\Lambda(P)$ for all $k$.
By \Cref{thm:slack-vi-upper}, $\Delta_{2k}(x,y)\le u_k(x)-l_k(y)$ with
$u_k=T^{k}\mathbf1$ and $l_k=T^{k}\mathbf0$, and $u_k,l_k\to w^\ast$ by
\Cref{thm:contraction}. Let $k\to\infty$. At an optimal successor the
right-hand side is $0$.
\end{proof}

\begin{lemma}[What a lifted certificate has to check]\label{ob:reduction}
Let $G$ be a stopping SSG whose non-sink vertices outside a set
$C\subseteq\Vmax$ are all average, let $\theta\colon V\to\mathbb R^{C}_{\ge0}$
be the \emph{lift} --- the unique solution of
$\theta_x=\tfrac12(\theta_{x^{(0)}}+\theta_{x^{(1)}})$ at average $x$,
$\theta_v=e_v$ on $C$, $\theta_{t_0}=\theta_{t_1}=0$ --- and let
$\Xi_{K}\subseteq\cdots\subseteq\Xi_0\subseteq\mathbb R^{C}_{\ge0}$ be
nonempty compact convex, with
$P_k:=\{x\mapsto w^\ast(x)+\langle\theta_x,\xi\rangle:\xi\in\Xi_k\}$.
Suppose
\begin{itemize}[leftmargin=2.6em]
\item[$(\alpha)$] $P_k\subseteq Q(G;L_Z,U_Z)$ for every $k\le K$;
\item[$(\beta)$] for every $v\in C$, every $y\in V$ and every $k<K$,
\[
  \max_{i}\Lambda(P_k)\bigl(v^{(i)},y\bigr)\ \ge\ \Lambda(P_{k+1})(v,y)
  \qquad\text{and}\qquad
  \max_{i}\Lambda_\times(P_k)\bigl(v^{(i)},y\bigr)\ \ge\
  \Lambda_\times(P_{k+1})(v,y).
\]
\end{itemize}
Then $(P_k)_{k\le K}$ satisfies hypotheses \textup{(i)} and \textup{(ii)}
of \Cref{ob:invariance}.
\end{lemma}

\begin{proof}
Write $\Lambda_k:=\Lambda(P_k)$, $\Lambda^{\times}_k:=\Lambda_\times(P_k)$.

\emph{Average vertices.} If $x$ is average then $w^\ast(x)$ and
$\theta_x$ are the means of those of $x^{(0)},x^{(1)}$, so
$z(x)=\tfrac12(z(x^{(0)})+z(x^{(1)}))$ for every $z\in P_{k+1}$; taking
$z$ attaining $\Lambda_{k+1}(x,y)$,
$\Lambda_{k+1}(x,y)=\tfrac12\sum_i(z(x^{(i)})-z(y))
\le\tfrac12\sum_i\Lambda_{k+1}(x^{(i)},y)\le\tfrac12\sum_i\Lambda_k(x^{(i)},y)$,
the (up) clause. The (down) clause at an average $y$ and the two
average--average composites are the same computation applied to
$\theta_x-\theta_y=\tfrac12\sum_i(\theta_{x^{(i)}}-\theta_{y^{(\pi(i))}})$.
Multiplicatively, for $z$ attaining $\Lambda^\times_{k+1}(x,y)$: the (up)
clause at an average $x$ is
$\tfrac12\sum_i\Lambda^\times_k(x^{(i)},y)\ge
\tfrac12\sum_iz(x^{(i)})/z(y)=z(x)/z(y)$; the (down) clause at an average
$y$ is the harmonic mean, which is monotone, so
$\mathrm{hm}\bigl(\Lambda^\times_k(x,y^{(0)}),\Lambda^\times_k(x,y^{(1)})\bigr)
\ge\mathrm{hm}\bigl(z(x)/z(y^{(0)}),z(x)/z(y^{(1)})\bigr)
=z(x)\big/\tfrac12\textstyle\sum_jz(y^{(j)})=z(x)/z(y)$; and the
average--average composite is the mediant inequality
$z(x)/z(y)=\sum_iz(x^{(i)})\big/\sum_iz(y^{(\pi(i))})
\le\max_iz(x^{(i)})/z(y^{(\pi(i))})$.

\emph{(down) at $v\in C$.} $P_m\subseteq Q(G)$ gives $z(v)\ge z(v^{(j)})$
pointwise, hence $\Lambda_m(x,v)\le\Lambda_m(x,v^{(j)})$ and
$\Lambda^\times_m(x,v)\le\Lambda^\times_m(x,v^{(j)})$ for all $m,x,j$;
with $\Lambda_{k+1}\le\Lambda_k$,
$\min_j\Lambda_k(x,v^{(j)})\ge\Lambda_k(x,v)\ge\Lambda_{k+1}(x,v)$, which
is the clause, and likewise multiplicatively. There are no $\Vmin$
vertices, so no dual case arises.

\emph{Base clauses.} By $(\alpha)$ every $z\in P_{k+1}$ has $z(t_0)=0$,
$z(t_1)=1$, $0\le z\le1$, $z\equiv0$ on $Z_0$ and $z\equiv1$ on $Z_1$.
Hence $\Lambda_{k+1}\le1$; $\Lambda_{k+1}(x,x)=0$;
$\Lambda_{k+1}(x,y)\le0$ when $x\in Z_0$ or $y\in Z_1$;
$\Lambda_{k+1}(t_0,t_1)=-1$; and $\Lambda^\times_{k+1}(x,y)=0$ when
$x\in Z_0$, $\Lambda^\times_{k+1}(x,y)\le1$ when $y\in Z_1$,
$\Lambda^\times_{k+1}(x,x)=1$.

\emph{Composites at two vertices of $C$.} By $(\beta)$ and the pointwise
fact just used,
$\Lambda_{k+1}(x,y)\le\max_i\Lambda_k(x^{(i)},y)
\le\max_i\min_j\Lambda_k(x^{(i)},y^{(j)})$ and
$\Lambda_{k+1}(x,y)\le\min_j\Lambda_{k+1}(x,y^{(j)})
\le\min_j\max_i\Lambda_k(x^{(i)},y^{(j)})$; identically in the
multiplicative register.

Every clause of \Cref{def:slack} and of \Cref{def:ratio} is now
accounted for.
\end{proof}

\subsection{The branch-compensated core}

\begin{definition}[$\mathrm{BC}(e,s)$]\label{ob:bc}
For $e\ge2$, $s\ge5$ put $u:=2^{-e}$, $\gamma:=2^{-s}$. Let
$\mathrm{CQ},\mathrm{CE},\mathrm{CH},\mathrm{Q}_0$ be the dyadic sink
chains of \Cref{lem:dyadic-row} realising the constants
$\tfrac14,\tfrac38,\tfrac34,\tfrac12+2\gamma$, of lengths $2,3,2,s-1$;
every vertex of them is average and reads only the two sinks.
$\mathrm{BC}(e,s)$ has $\Vmax=\{v_0,v_1,v_2\}$, $\Vmin=\emptyset$, all
other non-sinks average, and, with $o:=3-i$,
$\mathrm{HI}_1:=t_1$, $\mathrm{HI}_2:=\mathrm{CH}$,
$\mathrm{LO}_1:=\mathrm{CQ}$, $\mathrm{LO}_2:=\mathrm{CE}$:
\[
\begin{aligned}
&\alpha^{i}_q\to(v_{o},\alpha^{i}_{q+1})\ (q<e), &&\alpha^{i}_{e}\to(v_{o},\mathrm{HI}_i),\\
&A^{i}_1\to(\alpha^{i}_1,A^{i}_2),&&
A^{i}_q\to(v_{o},A^{i}_{q+1})\ (2\le q<e),\qquad A^{i}_{e}\to(v_{o},\mathrm{LO}_i),\\
&\beta^{i}_q\to(v_{i},\beta^{i}_{q+1})\ (q<2e),&&\beta^{i}_{2e}\to(v_{i},\mathrm{HI}_i),\\
&B^{i}_1\to(\beta^{i}_1,B^{i}_2),&&
B^{i}_q\to(v_{i},B^{i}_{q+1})\ (2\le q<2e),\qquad B^{i}_{2e}\to(v_{i},\mathrm{LO}_i),\\
&v_i\to(A^{i}_1,B^{i}_1)\ (i=1,2),&&\\
&g\to(g,t_1),\qquad P\to(g,t_0),&&\\
&D_q\to(v_1,D_{q+1})\ (q<s-3),&&D_{s-3}\to(v_1,t_0),\\
&Q\to(\mathrm{Q}_0,D_1),\qquad v_0\to(Q,P),&&
\end{aligned}
\]
with start vertex $v_0$.
\end{definition}

\begin{proposition}[Basic data]\label{ob:bc-values}
$\mathrm{BC}(e,s)$ is a stopping one-player SSG with $N=12e+2s+11$.
Put $\eta_1:=\tfrac12$, $\eta_2:=\tfrac14$ (so
$\eta_i=w^\ast(\mathrm{HI}_i)-\tfrac12$) and $\zeta_i:=\tfrac12\eta_i$
(so $\zeta_i=\tfrac12-w^\ast(\mathrm{LO}_i)$). Then
\[
\begin{aligned}
&w^\ast(v_0)=w^\ast(v_1)=w^\ast(v_2)=w^\ast(P)=w^\ast(A^{i}_1)=w^\ast(B^{i}_1)=\tfrac12,
\qquad w^\ast(g)=1,\qquad w^\ast(Q)=\tfrac12-\gamma,\\
&w^\ast(\alpha^{i}_q)=\tfrac12+2^{-(e-q+1)}\eta_i,\qquad
 w^\ast(A^{i}_q)=\tfrac12-2^{-(e-q+1)}\zeta_i\quad(q\ge2),\\
&w^\ast(\beta^{i}_q)=\tfrac12+2^{-(2e-q+1)}\eta_i,\qquad
 w^\ast(B^{i}_q)=\tfrac12-2^{-(2e-q+1)}\zeta_i\quad(q\ge2),\qquad
 w^\ast(D_q)=\tfrac12-2^{-(s-1-q)} .
\end{aligned}
\]
Hence $Z_0=\{t_0\}$, $Z_1=\{g,t_1\}$, and $v_0$ is the \emph{unique}
value-distinguishing controlled vertex, with $w^\ast(P)-w^\ast(Q)=\gamma$
and $P$ optimal. The lift for $C=\{v_0,v_1,v_2\}$, in coordinates
$(e_0,e_1,e_2)$, is $\theta_{v_j}=e_j$, $\theta=0$ at $g$, $P$, the two
sinks and every dyadic chain, and
\[
\begin{aligned}
&\theta_{\alpha^{i}_q}=\bigl(1-2^{-(e-q+1)}\bigr)e_{o}\ (q\ge1),\quad
 \theta_{A^{i}_q}=\bigl(1-2^{-(e-q+1)}\bigr)e_{o}\ (q\ge2),
&&\theta_{A^{i}_1}=a\,e_{o},\\
&\theta_{\beta^{i}_q}=\bigl(1-2^{-(2e-q+1)}\bigr)e_{i}\ (q\ge1),\quad
 \theta_{B^{i}_q}=\bigl(1-2^{-(2e-q+1)}\bigr)e_{i}\ (q\ge2),
&&\theta_{B^{i}_1}=b\,e_{i},\\
&\theta_{D_q}=\bigl(1-2^{-(s-2-q)}\bigr)e_1,&&\theta_Q=c\,e_1,
\end{aligned}
\qquad
a:=1-\tfrac32u,\quad b:=1-\tfrac32u^{2},\quad c:=\tfrac12-4\gamma .
\]
Every $\theta_x$ is $0$, $e_0$, or $c'e_1$ or $c'e_2$ with
$c'\in(0,1]$; every $y$ with $\theta_y=c'e_2$ has
$w^\ast(y)\ge\tfrac14+\tfrac{c'}4$, and every $y$ with $\theta_y=c'e_1$
has $w^\ast(y)\ge\tfrac{c'}2$.
\end{proposition}

\begin{proof}
Counting: $3+3+\bigl(2+3+2+(s-1)\bigr)+2\cdot6e+(s-3)=12e+2s+9$
non-sinks. Stopping (\Cref{lem:trapchar}): every dyadic chain vertex has
a sink successor, so lies outside any trap $U$; then
$\alpha^{i}_e,\beta^{i}_{2e}\notin U$, inductively $\alpha^{i}_1,
\beta^{i}_1\notin U$, hence the average vertices $A^{i}_1,B^{i}_1\notin U$,
hence $v_1,v_2\notin U$ (no successor inside); then $D_{s-3}$, every
$D_q$, $Q$, $g$, $P$ and $v_0$. So $U=\emptyset$.

The displayed vector is a fixed point of $T$. On the dyadic chains this is
\Cref{lem:dyadic-row}; $w^\ast(g)=\tfrac12(1+1)$ and $w^\ast(P)=\tfrac12(1+0)$.
Along $\alpha^{i}$, writing $y_q:=w^\ast(\alpha^{i}_q)-\tfrac12$, the last
vertex gives $y_e=\tfrac12\bigl(w^\ast(\mathrm{HI}_i)-\tfrac12\bigr)=\tfrac12\eta_i$
and $y_q=\tfrac12y_{q+1}$, which is the display; the same computation on
$\beta^{i}$, and on $A^{i}$, $B^{i}$ for $q\ge2$ with
$-\tfrac12\zeta_i$ in place of $\tfrac12\eta_i$. Then
\[
  w^\ast(A^{i}_1)=\tfrac12\bigl(w^\ast(\alpha^{i}_1)+w^\ast(A^{i}_2)\bigr)
  =\tfrac12+\tfrac12\bigl(2^{-e}\eta_i-2^{-(e-1)}\zeta_i\bigr)=\tfrac12
\]
because $\eta_i=2\zeta_i$, and identically $w^\ast(B^{i}_1)=\tfrac12$;
so $w^\ast(v_i)=\max(\tfrac12,\tfrac12)=\tfrac12$. Backwards from
$w^\ast(D_{s-3})=\tfrac12(\tfrac12+0)=\tfrac14$ gives the $D$-chain,
whence $w^\ast(D_1)=\tfrac12-2^{-(s-2)}=\tfrac12-4\gamma$ and
$w^\ast(Q)=\tfrac12\bigl((\tfrac12+2\gamma)+(\tfrac12-4\gamma)\bigr)
=\tfrac12-\gamma$, and $w^\ast(v_0)=\max(\tfrac12-\gamma,\tfrac12)=\tfrac12$.
Since $G$ is stopping, \Cref{cor:comparison} identifies this fixed point
with $w^\ast$. The qualitative sets are read off: only $t_0$ has value
$0$, only $g$ and $t_1$ have value $1$. $v_1$ and $v_2$ are not
value-distinguishing; $v_0$ is, with gap $\gamma$.

The lift is the same backward substitution with $\theta_{v_j}=e_j$; a
chain of $r$ average vertices reading $v_j$ and ending at a
$\theta$-free vertex carries $\theta=(1-2^{-r})e_j$ at its head, and
$\theta_{A^{i}_1}=\tfrac12\bigl((1-2^{-e})+(1-2^{-(e-1)})\bigr)e_o
=(1-\tfrac32u)e_o$, $\theta_{B^{i}_1}=(1-\tfrac32u^{2})e_i$,
$\theta_Q=\tfrac12\bigl(0+(1-2^{-(s-3)})\bigr)e_1=(\tfrac12-4\gamma)e_1$.

Last sentence. A vertex $y$ with $\theta_y=c'e_2$, $c'>0$, is one of
$v_2$, $\alpha^{1}_\bullet$, $A^{1}_\bullet$, $\beta^{2}_\bullet$,
$B^{2}_\bullet$; its value is $\tfrac{c'}2+(1-c')\tau_y$ where $\tau_y$
is a convex combination of the terminal values it reaches, namely
$1,\tfrac34,\tfrac14,\tfrac38$, all $\ge\tfrac14$, so
$w^\ast(y)\ge\tfrac{c'}2+\tfrac{1-c'}4=\tfrac14+\tfrac{c'}4$. A vertex
with $\theta_y=c'e_1$ may in addition reach the terminal $t_0$ (along
$D$), so only $w^\ast(y)\ge\tfrac{c'}2$ is available --- and it is
attained, $w^\ast(D_q)=\tfrac12\theta_{D_q}$ exactly.
\end{proof}

\begin{proposition}[No sound local matching calculus decides
$v_0$]\label{ob:bc-lmc}
In $\mathrm{BC}(e,s)$ the pair $(Q,P)$ carries the frame of the pair
$(4,1)$ of $G_8$ (\Cref{prop:simorder-stalls,lem:fooling-partner}) and the
pair $(v_0,P)$ carries the frame of $(5,1)$, in each case with the same
set of value-true atoms. Hence no sound local matching calculus
(\Cref{def:lmc}) derives $(Q,P)$ or $(v_0,P)$, in either reading; and
$(v_0,Q)$, $(P,Q)$ are value-false. So no sound local matching calculus
decides $v_0$ by any of the four tests of \Cref{rem:own-successor}
available to an order-only mechanism. Moreover the preorder $\preceq$ of
\Cref{def:simorder} decides neither $v_1$ nor $v_2$.
\end{proposition}

\begin{proof}
The frames. Writing $x,x^{(0)},x^{(1)},y,y^{(0)},y^{(1)}$ for the six
symbols, $(Q,P)$ interprets them at $Q,\mathrm{Q}_0,D_1,P,g,t_0$, of types
$(\mathrm{avg},\mathrm{avg},\mathrm{avg},\mathrm{avg},\mathrm{avg},t_0)$,
all six distinct, with value tuple
$(\tfrac12-\gamma,\tfrac12+2\gamma,\tfrac12-4\gamma;\tfrac12,1,0)$; the
value-true atoms are exactly
$\Pi=\{(x,y^{(0)}),(x^{(0)},y^{(0)}),(x^{(1)},y),(x^{(1)},y^{(0)})\}$,
which is what \Cref{lem:fooling-partner} computes at $(4,1)$ of $G_8$
(value tuple $(\tfrac25,\tfrac35,\tfrac15;\tfrac12,1,0)$) and at
$(a_0,d)$ of $G_{\mathrm f}$. The proof of \Cref{thm:matching-barrier}
applies verbatim: a sound clause deriving $(Q,P)$ is contained in $\Pi$,
and the same clause derives $(a_0,d)$ in $G_{\mathrm f}$, where
$w^\ast(a_0)>w^\ast(d)$.

$(v_0,P)$ interprets the symbols at $v_0,Q,P,P,g,t_0$, of types
$(\max,\mathrm{avg},\mathrm{avg},\mathrm{avg},\mathrm{avg},t_0)$, with
equality pattern $\{x^{(1)}=y\}$ and value-true atoms
$\Pi\cup\{(x^{(0)},y)\}$, the extra atom interpreting as $(Q,P)$. That is
precisely the frame, equality pattern and true-atom set of $(5,1)$ in
$G_8$. Since $(Q,P)$ is underivable, a clause deriving $(v_0,P)$ is
contained in $\Pi$, and the game $K_\sharp$ of \Cref{thm:matching-barrier}
fools every such clause.

For $\preceq$ at $v_1$ (and symmetrically $v_2$) put
$\mathcal S_0:=\mathcal V:=\{(x,y):w^\ast(x)\le w^\ast(y)\}$ and
$\mathcal S_{j+1}:=\mathcal V\cap\mathcal F(\mathcal S_j)$; since
$\preceq\subseteq\mathcal V$ (\Cref{thm:simorder-sound}) and
$\preceq=\mathcal F(\preceq)$, induction gives
$\preceq\subseteq\mathcal S_j$ for every $j$. Write
$\rho:=2^{-e}\eta_1$ and $\rho':=2^{-2e}\eta_1$, so that the successors of
$A^{1}_1$ have values $\tfrac12\pm\rho$ and those of $B^{1}_1$ have values
$\tfrac12\pm\rho'$, with $0<\rho'<\rho$. Then
$(B^{1}_1,A^{1}_1)\notin\mathcal S_1$: the (down) clause at $A^{1}_1$
needs $(B^{1}_1,A^{1}_2)$, i.e.\ $\tfrac12\le\tfrac12-\rho$; the (up)
clause at $B^{1}_1$ needs $(\beta^{1}_1,A^{1}_1)$, i.e.\
$\tfrac12+\rho'\le\tfrac12$; the (match) clause needs, for
$\pi=\mathrm{id}$, $(B^{1}_2,A^{1}_2)$, i.e.\
$\tfrac12-\rho'\le\tfrac12-\rho$, false since $\rho'<\rho$, and for the
transposition $(\beta^{1}_1,A^{1}_2)$, comparing a value above $\tfrac12$
with one below. Symmetrically $(A^{1}_1,B^{1}_1)\notin\mathcal S_1$, the
$\pi=\mathrm{id}$ branch failing at $(\alpha^{1}_1,\beta^{1}_1)$, i.e.\ at
$\tfrac12+\rho\le\tfrac12+\rho'$. Hence
$(v_1,A^{1}_1)\notin\mathcal S_2$: its (up) clause needs
$(B^{1}_1,A^{1}_1)\in\mathcal S_1$, its (down) clause needs
$(v_1,A^{1}_2)\in\mathcal V$, i.e.\ $\tfrac12\le\tfrac12-\rho$, and no
match clause applies, the types differing; likewise
$(v_1,B^{1}_1)\notin\mathcal S_2$. The remaining two tests at $v_1$ are
the pair tests, which are the two pairs just excluded.

Two remarks on the shape of the construction. First, the two halves of the
engine must carry \emph{different} straddle amplitudes: if
$\eta_1=\eta_2$ the match clause relates $A^{1}_1$ with $A^{2}_1$ and
$B^{1}_1$ with $B^{2}_1$ in both directions, and then
$v_1\preceq v_2\preceq v_1$, so $\preceq$ contributes the cuts
$z(v_1)=z(v_2)$; those cuts collapse $\Xi_k$ of \Cref{ob:bc-cert} to its
diagonal, and \Cref{ob:bc-lower}'s clause on seeding by $M1$ fails.
(Verified: with $\eta_1=\eta_2$ the preorder still decides no controlled
vertex, but its cuts are violated by the certificate at every round.)
Second, $\preceq$ is already transitive on $\mathrm{BC}(e,s)$ --- verified
exactly at the ten sizes of \Cref{ob:measured} --- so reading
\Cref{def:simorder} with transitive closure changes nothing. For a
\emph{general} sound local matching calculus read with transitive closure
we make no claim: \Cref{thm:matching-barrier}'s closure paragraph turns on
the value multiset of $G_8$, which has only three vertices in the relevant
interval, and $\mathrm{BC}(e,s)$ has many.
\end{proof}

\begin{theorem}[The certificate]\label{ob:bc-cert}
With $a,b,c$ as in \Cref{ob:bc-values} put $\ell:=\gamma/c$,
$M_k:=\tfrac14a^{k}$,
\[
  \Xi_k:=\bigl\{\xi\in\mathbb R^{3}:\ \ell\le\xi_1,\xi_2\le M_k,\
  \xi_1\ge a\xi_2,\ \xi_2\ge a\xi_1,\ \xi_0=c\xi_1-\gamma\bigr\},\qquad
  P_k:=\{x\mapsto w^\ast(x)+\langle\theta_x,\xi\rangle:\xi\in\Xi_k\},
\]
and let $K$ be the largest $k$ with $aM_k\ge\ell$. Then $(P_k)_{k\le K}$
satisfies the hypotheses of \Cref{ob:invariance}.
\end{theorem}

\begin{proof}
Write $h_k(\cdot):=\max_{\Xi_k}\langle\cdot,\xi\rangle$, so
$\Lambda(P_k)(x,y)=w^\ast(x)-w^\ast(y)+h_k(\theta_x-\theta_y)$; on $\Xi_k$,
$\langle e_0,\xi\rangle=c\xi_1-\gamma$. For $k\le K$, $\Xi_k$ is a
hexagon with vertices
$(\ell,\ell),(\ell,\ell/a),(\ell/a,\ell),(M_k,M_k),(M_k,aM_k),(aM_k,M_k)$
in the coordinates $(\xi_1,\xi_2)$, all admissible because $aM_k\ge\ell$
(hence $\ell/a\le M_k$) and $a<1$; and $\Xi_{k+1}\subseteq\Xi_k$ because
$M_{k+1}=aM_k$. We verify $(\alpha)$ and $(\beta)$ of
\Cref{ob:reduction}.

$(\alpha)$. The average rows hold by definition of $\theta$. On $\Xi_k$,
$\xi_1,\xi_2\ge\ell>0$ and $\xi_0=c\xi_1-\gamma\ge0$, so $\xi\ge0$ and,
$\theta$ being nonnegative, $z\ge w^\ast\ge0$. Every $\theta_x$ is a
convex combination of $0$ and the $e_j$, and
$\langle e_0,\xi\rangle\le cM_k\le M_k$, so
$\langle\theta_x,\xi\rangle\le M_k\le\tfrac14$; the largest value carried
by a vertex with $\theta_x\ne0$ is $\tfrac34$ (at $\alpha^{i}_e$ and
$\beta^{i}_{2e}$), so $z\le1$; at $g$ and $t_1$, $\theta=0$ and $z=1$; at
$t_0$, $z=0$: these are the box rows and the $Z$ pins. Max rows: at
$v_1$, $z(v_1)\ge z(A^{1}_1)$ reads $\xi_1\ge a\xi_2$ and
$z(v_1)\ge z(B^{1}_1)$ reads $\xi_1\ge b\xi_1$, true because $\xi_1\ge0$
and $b\le1$; at $v_2$ symmetrically; at $v_0$, $z(v_0)\ge z(P)$ reads
$\xi_0\ge0$ and $z(v_0)\ge z(Q)$ reads $\xi_0\ge c\xi_1-\gamma$, an
equality.

$(\beta)$ at $v_0$. On every $P_k$ one has
$z(v_0)=\tfrac12+\xi_0=\tfrac12-\gamma+c\xi_1=z(Q)$ \emph{identically},
so $\Lambda(P_{k+1})(v_0,y)=\Lambda(P_{k+1})(Q,y)\le\Lambda(P_k)(Q,y)$
and $\Lambda_\times(P_{k+1})(v_0,y)=\Lambda_\times(P_{k+1})(Q,y)
\le\Lambda_\times(P_k)(Q,y)$, and $Q$ is one of the two successors.

$(\beta)$ at $v_1$ (the case of $v_2$ exchanges $e_1$ and $e_2$). Since
$w^\ast(A^{1}_1)=w^\ast(B^{1}_1)=w^\ast(v_1)=\tfrac12$, the additive
condition is
\[
  \max\bigl(h_k(a\,e_2-\theta_y),\ h_k(b\,e_1-\theta_y)\bigr)\ \ge\
  h_{k+1}(e_1-\theta_y)\qquad\text{for every }y,
\]
and by \Cref{ob:bc-values} there are four cases.
\emph{$\theta_y=0$}: the left side is at least $bM_k$, the right side is
$M_{k+1}=aM_k$, and $b>a$.
\emph{$\theta_y=c'e_1$}: the right side is $(1-c')M_{k+1}=a(1-c')M_k$,
and at $\xi=(aM_k,M_k)\in\Xi_k$ the first term is
$aM_k-c'aM_k=a(1-c')M_k$.
\emph{$\theta_y=c'e_2$}: as $\xi_2\ge a\xi_1$, the right side is at most
$(1-c'a)M_{k+1}=a(1-c'a)M_k$; at $\xi=(M_k,aM_k)\in\Xi_k$ the second term
is $(b-c'a)M_k$; and
$(b-c'a)-a(1-c'a)=(b-a)-c'a(1-a)=\tfrac32u(1-u)-\tfrac32u\,c'a
=\tfrac32u\,(1-u-c'a)\ge0$ because $c'a\le a=1-\tfrac32u\le1-u$.
\emph{$\theta_y=e_0$}: the right side is $\gamma+(1-c)M_{k+1}$, and at
$\xi=(aM_k,M_k)$ the first term is $aM_k-caM_k+\gamma=\gamma+(1-c)M_{k+1}$.

The multiplicative condition at $v_1$ is
\[
  \max\Bigl(\max_{\Xi_k}\frac{\tfrac12+a\xi_2}{z_\xi(y)},\
  \max_{\Xi_k}\frac{\tfrac12+b\xi_1}{z_\xi(y)}\Bigr)\ \ge\
  \max_{\Xi_{k+1}}\frac{\tfrac12+\xi_1}{z_\xi(y)},\qquad
  z_\xi(y)=w^\ast(y)+\langle\theta_y,\xi\rangle .
\]
The denominator vanishes only at $y=t_0$, where both sides are $+\infty$.
Elsewhere all three maxima are attained at vertices of the hexagons, a
linear-fractional function with positive denominator being both
quasi-convex and quasi-concave.
\emph{$\theta_y=0$}: the denominator is constant and the claim is
$bM_k\ge aM_k$.
\emph{$\theta_y=c'e_1$}: both sides depend only on $\xi_1$, through
$t\mapsto(\tfrac12+t)/(w^\ast(y)+c't)$, which is nondecreasing because
$w^\ast(y)\ge c'/2$ (\Cref{ob:bc-values}); so the right side is
$(\tfrac12+M_{k+1})/(w^\ast(y)+c'M_{k+1})$, and the first term at
$\xi=(aM_k,M_k)\in\Xi_k$ is the same number.
\emph{$\theta_y=c'e_2$}: the right-hand function increases in $\xi_1$ and
decreases in $\xi_2$, so its maximum lies on the edge $\xi_2=a\xi_1$ of
$\Xi_{k+1}$, where it is $g(t)=(\tfrac12+t)/(w^\ast(y)+c'at)$,
$t\in[\ell/a,M_{k+1}]$; $g$ is nondecreasing because
$w^\ast(y)\ge\tfrac14+\tfrac{c'}4\ge\tfrac{c'}2\ge\tfrac{c'a}2$, so the
maximum is $g(M_{k+1})=(\tfrac12+aM_k)/(w^\ast(y)+c'a^{2}M_k)$. The
second term at $\xi=(M_k,aM_k)\in\Xi_k$ is
$(\tfrac12+bM_k)/(w^\ast(y)+c'aM_k)$; writing $W:=w^\ast(y)$, $M:=M_k$,
the required inequality
$(\tfrac12+bM)(W+c'a^{2}M)\ge(\tfrac12+aM)(W+c'aM)$ expands, after
dividing by $\tfrac32u$, to $(1-u)W\ge c'a(\tfrac12+uaM)$. Since
$M\le\tfrac14$ and $a\le1$ the right side is at most
$c'(1-\tfrac32u)(\tfrac12+\tfrac u4)=c'(\tfrac12-\tfrac u2-\tfrac38u^{2})$,
while $W\ge\tfrac14+\tfrac{c'}4$ makes the left side at least
$\tfrac14(1+c')(1-u)$, and
\[
  \tfrac14(1+c')(1-u)-c'\bigl(\tfrac12-\tfrac u2-\tfrac38u^{2}\bigr)
  \;=\;\tfrac{(1-c')(1-u)}4+\tfrac{3c'u^{2}}8\;\ge\;0 .
\]
\emph{$\theta_y=e_0$}: $z_\xi(v_0)=\tfrac12-\gamma+c\xi_1$ and
$t\mapsto(\tfrac12+t)/(\tfrac12-\gamma+ct)$ increases because
$\tfrac12-\gamma\ge\tfrac c2=\tfrac14-2\gamma$; the right side is
attained at $\xi_1=M_{k+1}=aM_k$, and equals the first term at
$\xi=(aM_k,M_k)\in\Xi_k$.

This is $(\beta)$; \Cref{ob:reduction} finishes the proof. The bound
$M_0\le\tfrac14$ is used exactly twice: in the box row of $(\alpha)$ and
in the third multiplicative case.
\end{proof}

\begin{theorem}[The lower bound]\label{ob:bc-lower}
For every $k\le K$,
\[
  \Lambda(P_k)(v_0,P)=\Lambda(P_k)(Q,P)=cM_k-\gamma>0,\qquad
  \Lambda(P_k)(Q,v_0)=0,
\]
\[
  \Lambda(P_k)(v_1,A^{1}_1)=(1-a^{2})M_k>0,\quad
  \Lambda(P_k)(v_1,B^{1}_1)=(1-b)M_k>0,\quad
  \Lambda(P_k)(A^{1}_1,B^{1}_1)=a(1-b)M_k>0,\quad
  \Lambda(P_k)(B^{1}_1,A^{1}_1)=(b-a^{2})M_k>0,
\]
symmetrically at $v_2$; each of the six strictly positive quantities is
$>1$ in the multiplicative register, and $\Lambda_\times(P_k)(Q,v_0)=1$
exactly. Consequently no member of $\mathfrak D_0$ ---
in particular none of $M2$, $M2T$, $M4$, $M5$, $M6$, the $Z$-seeded
own-successor hybrid, nor any combination or mutual seeding of them ---
decides any controlled vertex of $\mathrm{BC}(e,s)$ within
\[
  K+1\;=\;\Bigl\lfloor\frac{\ln\bigl(c/4\gamma\bigr)}{\ln(1/a)}\Bigr\rfloor
  \;\ge\;\Bigl\lfloor\frac{(2^{e+1}-3)(s-4)\ln2}{3}\Bigr\rfloor
\]
operator applications; with $s$ fixed and $e=(N-2s-11)/12$ that is
$2^{\Omega(N)}$. Adjoining \Cref{def:simorder} changes nothing: every
point of every $P_k$ satisfies every cut $z(x)\le z(y)$ with
$x\preceq y$, so the composite that runs $M1$ first and seeds the rest
with its verdicts obeys the same bound, and by \Cref{ob:bc-lmc} $M1$
itself decides nothing.
\end{theorem}

\begin{proof}
The displayed quantities. $\Lambda(P_k)(v_0,P)=h_k(e_0)=cM_k-\gamma$,
which is positive for $k\le K$ because $M_k\ge aM_k\ge\ell=\gamma/c$ and
$a<1$; $\Lambda(P_k)(Q,P)=-\gamma+h_k(c\,e_1)=cM_k-\gamma$;
$\Lambda(P_k)(Q,v_0)=-\gamma+h_k(c\,e_1-e_0)=-\gamma+\gamma=0$, since
$c\xi_1-\xi_0\equiv\gamma$ on $\Xi_k$;
$\Lambda(P_k)(v_1,A^{1}_1)=h_k(e_1-a\,e_2)=(1-a^{2})M_k$, attained at
$(M_k,aM_k)$; $\Lambda(P_k)(v_1,B^{1}_1)=h_k((1-b)e_1)=(1-b)M_k$;
$\Lambda(P_k)(A^{1}_1,B^{1}_1)=h_k(a\,e_2-b\,e_1)=a(1-b)M_k$, attained at
$(aM_k,M_k)$; and $\Lambda(P_k)(B^{1}_1,A^{1}_1)=h_k(b\,e_1-a\,e_2)
=(b-a^{2})M_k$, attained at $(M_k,aM_k)$ and positive because
$b-a^{2}=3u(1-\tfrac54u)>0$. Each of these being $>0$, \Cref{ob:sign}
makes the corresponding multiplicative quantity $>1$; and
$\Lambda_\times(P_k)(Q,v_0)=1$ because $z(Q)=z(v_0)>0$ identically on
$P_k$.

Now let $\mathcal A\in\mathfrak D_0$ report after $k\le K$ applications
from a sound state. At $v_0$: clause~(i) needs $\Delta(v_0,P)\le0$ or
$\Delta(v_0,Q)\le0$, but $\Delta(v_0,P)\ge\Lambda(P_k)(v_0,P)>0$ by
\Cref{ob:invariance} and $\Delta(v_0,Q)\ge w^\ast(v_0)-w^\ast(Q)=\gamma>0$
by soundness; clause~(ii) needs $\Delta(Q,v_0)<0$ or $\Delta(P,v_0)<0$,
but $\Delta(Q,v_0)\ge\Lambda(P_k)(Q,v_0)=0$ and
$\Delta(P,v_0)\ge w^\ast(P)-w^\ast(v_0)=0$; the pair tests need
$\Delta(Q,P)\le0$ or $\Delta(P,Q)\le0$, but
$\Delta(Q,P)\ge\Lambda(P_k)(Q,P)>0$ and $\Delta(P,Q)\ge\gamma>0$. At
$v_1$ and $v_2$ all three of $w^\ast(v_i)$, $w^\ast(A^{i}_1)$,
$w^\ast(B^{i}_1)$ equal $\tfrac12$, so clause~(ii) is blocked by soundness
alone ($\Delta(A^{i}_1,v_i),\Delta(B^{i}_1,v_i)\ge0$), while clause~(i)
and \emph{both} pair tests --- which soundness does not block there, all
three values being equal --- are blocked by the four positive quantities
displayed. The multiplicative register is blocked entry by entry by
\Cref{ob:sign}, by $\Lambda_\times(P_k)(Q,v_0)=1$, and by
$w^\ast(v_0)/w^\ast(Q)=\tfrac12/(\tfrac12-\gamma)>1$.

The count. $K$ is the largest $k$ with $a\cdot\tfrac14a^{k}\ge\gamma/c$,
i.e.\ $a^{k+1}\ge4\gamma/c$, i.e.\ $k+1\le\ln(c/4\gamma)/\ln(1/a)$; and
$\ln(1/a)=-\ln(1-\tfrac32u)\le\tfrac{3u/2}{1-3u/2}=\tfrac{3}{2^{e+1}-3}$
while $c/4\gamma=2^{s-3}-1\ge2^{s-4}$ for $s\ge5$.
\end{proof}

\begin{corollary}[The edge]\label{ob:edge}
On $\mathrm{BC}(e,s)$ every $\sigma$ with $\sigma(v_0)=P$ has no strictly
switchable vertex, so $\val_\sigma=w^\ast$ by \Cref{lem:local-global};
all-switches therefore halts within one round from every $\sigma$. Hence
one application of the operator \textup{(E)} of \Cref{ob:mem} --- one
exact policy evaluation --- yields $\Delta=w^\ast(x)-w^\ast(y)$ and
decides every controlled vertex, while by \Cref{ob:bc-lower} no sequence
of $2^{o(N)}$ operators of $\mathcal R$ decides any. So \textup{(E)} is an
extension of $\mathfrak D$ that breaks the barrier, it is not in
$\mathfrak D$, and $\mathrm{BC}(e,s)$ separates $\mathfrak D$ from
$\mathfrak D^{+}$ by an exponential factor in cost.
\end{corollary}

\begin{proof}
Under such a $\sigma$ the two successors of $v_1$ and of $v_2$ have equal
value and $P$ is optimal at $v_0$, so $\val_\sigma=w^\ast$ and $S_\sigma$
is empty; from a $\sigma$ with $\sigma(v_0)=Q$ the unique strictly
switchable vertex is $v_0$, and one all-switches round reaches the
previous case. With $\Vmin=\emptyset$, $U=\val^{\tau}=w^\ast$ and
$L=\val_\sigma=w^\ast$, so \textup{(E)} returns $w^\ast(x)-w^\ast(y)$,
which is $-\gamma<0$ at $(Q,P)$.
\end{proof}

\begin{remark}[Measured]\label{ob:measured}
Everything above was verified in exact rational arithmetic, from the game,
at $(e,s)=(2,5),(3,5),(4,5),(5,5),(6,5),(2,6),(3,6),(2,7),(3,8),(4,8)$,
i.e.\ $N=45,57,69,81,93,47,59,49,63,75$. For each: the closed forms of
$w^\ast$ and $\theta$ (\Cref{ob:bc-values}) exactly; stopping by
\Cref{lem:trapchar}; $w^\ast$ independently by enumeration of the $2^{3}$
positional profiles; $Z_0=\{t_0\}$, $Z_1=\{g,t_1\}$; the two frame
identities of \Cref{ob:bc-lmc}; the greatest fixed point $\preceq$ of
\Cref{def:simorder} computed outright, found to be transitive, and to
relate none of the twelve pairs the four tests of
\Cref{rem:own-successor} read at $v_0,v_1,v_2$; and, at every $k\le K$,
$P_k\subseteq Q(G;L_Z,U_Z)$, every clause of \Cref{def:slack} at
$\Lambda(P_k)$ at least $\Lambda(P_{k+1})$ at all $N^{2}$ entries, every
clause of \Cref{def:ratio} at $\Lambda_\times(P_k)$ at least
$\Lambda_\times(P_{k+1})$ at all $N^{2}$ entries, every point of $P_k$
obeying every $\preceq$-cut, and none of the six additive and six
multiplicative tests of \Cref{rem:own-successor} (two for clause~(i), two
for clause~(ii), two pair tests) firing at any of $v_0,v_1,v_2$: zero
violations of any kind. The round counts $K+1$ are
$2,5,11,22,46$ at $s=5$ and $e=2,\dots,6$ --- the doubling of
\Cref{ob:bc-lower} --- and $4,9$ at $s=6$, $5$ at $s=7$, $16,34$ at
$s=8$; the same certificate check was run out to $e=7$ at $s=5$
($N=105$, $93$ rounds), $e=6$ at $s=8$ ($N=99$, $144$ rounds) and
$e=4$ at $s=10$ ($N=79$, $49$ rounds), again with no violation, and the
closed form $K+1=\lfloor\ln(c/4\gamma)/\ln(1/a)\rfloor$ agrees with the
measured count at every one of the fourteen sizes. The slack calculus,
its transitive closure and the ratio calculus were also run directly on
$\mathrm{BC}(2,5)$ for ten rounds, plain and seeded with $\preceq$:
silent throughout. A direct run of the full hybrid of \Cref{def:hybrid}
was attempted and cut short for time (one round on $N=45$ is some two
thousand exact simplex re-optimisations); the transport step is covered
by \Cref{ob:invariance}'s \textup{(D3)} case, which is where the
certificate does that work.

The scope of \Cref{ob:bc-lower} is $\mathfrak D_0$ together with the one
\textup{(D8)} step that produces $\preceq$. A \textup{(D8)} step over
some \emph{other} finite grade set is not covered: its output is a sound
matrix, but we have no argument that the cuts it supplies are respected by
$P_k$, and \Cref{ob:invariance}'s induction cannot start inside a
coinductive iteration, which begins unsound. What can be said in general
is the coinductive bound: any sound fixed point $\Delta_\infty$ of
$\rho_\Gamma\Theta$ satisfies
$\Delta_\infty\ge\rho_\Gamma\Theta(\rho_\Gamma D)$, since
$\Delta_\infty\ge D$ forces $\Delta_\infty\ge\rho_\Gamma D$ and
$\Theta$ is monotone; for $\Gamma=\{0,1\}$ and $\Theta=\mathcal L$ this
is $\preceq\subseteq\mathcal F(\mathcal V)$, the tool used in
\Cref{ob:bc-lmc}.
\end{remark}

\begin{remark}[What this covers that \Cref{prop:locality} and
\Cref{thm:seed-dichotomy} do not]\label{ob:relate}
\Cref{prop:locality} rules out rules reading a bounded neighbourhood;
every mechanism in $\mathfrak D$ reads the whole instance, and
\Cref{prop:simorder-stalls} already records that $M1$ beats every bounded
radius, so \Cref{ob:bc-lower} is not a restatement of it.
\Cref{thm:seed-dichotomy} says that a stalling instance of $M(p,k)$ is one
on which all-switches has not converged in $p$ rounds; on
$\mathrm{BC}(e,s)$ all-switches converges in one round, so the dichotomy
predicts --- correctly, \Cref{ob:edge} --- that $M(1,0)$ decides it. The
content of \Cref{ob:bc-lower} is that everything on the propagation side
of that dichotomy fails anyway, so the frontier the dichotomy locates is
not an artefact of any one mechanism but of the entire closure of them.

$\mathrm{BC}(e,s)$ is the first family here that defeats $M1$ and $M6$
together with the hybrid: $H_m$, $CC$, $WD$ and $CV$ are all decided at
round zero by $M1$, and $CV$ is decided by $M6$ within a few dozen rounds
(\Cref{prop:cv-measured}). What it gives up is the constant denominator:
$D(\mathrm{BC}(e,s))=2^{\Theta(e)}$, so unlike $CV$ it does not lie in the
class of \Cref{thm:few-denominator}. Like every family here it is easy:
being one-player, \Cref{thm:one-player} solves it by one linear
programme.
\end{remark}